\documentclass[../../main.tex]{subfiles}% 注意这里的文件路径不能用 ./main.tex ,否则用latexmk编译子文件会报错
\graphicspath{{\subfix{./image/}}{\subfix{../../image/}}} % 指定图片引用路径,后续可以直接使用图片文件名
% 如果图片存放在上级目录,则使用 ../../image/ 作为备用路径

% 例如:
% \begin{figure}[H]
% \centering
% \includegraphics[width=0.3\textwidth]{图.png}
% \caption{}
% \label{figure:图}
% \end{figure}
% 注意:上述\label{}一定要放在\caption{}之后,否则引用图片序号会只会显示??.

\begin{document}

\section{与核相联的容量}\label{section:ch03-核容量}

\subsection{预备知识}\label{subsection:ch03-核预备}

设 $g:\mathbb{R}^N\times\mathbb{R}^N\to[0,\infty]$ 为连续对称函数。

\begin{enumerate}[(i)]
\item 对任意 $\mu\in\mathscr{R}^+(\mathbb{R}^N)$ 与 $x\in\mathbb{R}^N$，置
\begin{align}
g(x,\mu)\triangleq \int_{\mathbb{R}^N}g(x,y)\,\mathrm{d}\mu(y).\label{align:ch03-3.21a}
\end{align}
\item 对任意 $\nu\in\mathscr{R}^+(\mathbb{R}^N)$ 与 $y\in\mathbb{R}^N$，置
\begin{align}
g(\nu,y)\triangleq \int_{\mathbb{R}^N}g(x,y)\,\mathrm{d}\nu(x).\label{align:ch03-3.21b}
\end{align}
\item 对任意 $\mu,\nu\in\mathscr{R}^+(\mathbb{R}^N)$，置
\begin{align}
g(\nu,\mu)\triangleq \int_{\mathbb{R}^N}g(x,\mu)\,\mathrm{d}\nu(x)
=\int_{\mathbb{R}^N}\int_{\mathbb{R}^N}g(x,y)\,\mathrm{d}\mu(y)\,\mathrm{d}\nu(x).\label{align:ch03-3.21c}
\end{align}
\end{enumerate}

由 Tonelli 定理，$g(\nu,\mu)=g(\mu,\nu)$。

设 $g(x,y)\equiv g(x-y)$（$x,y\in\mathbb{R}^N$），其中 $g:\mathbb{R}^N\to[0,\infty)$ 为径向递减连续函数。则由平移不变性，$g(x,\mu)=(g*\mu)(x)$，$g(\nu,y)=(g*\nu)(y)$，且
\begin{align}
g(\nu,\mu)=\int_{\mathbb{R}^N}(g*\mu)\,\mathrm{d}\nu.\label{align:ch03-3.22}
\end{align}

\begin{lemma}\label{lemma:ch03-下界估计}
设 $g(x,y)\equiv g(x-y)$，$g$ 径向递减。设 $\mu\in\mathscr{R}^+(\mathbb{R}^N)$，$E\subseteq\mathbb{R}^N$ 为 $\mu$-可测集。则
\begin{align}
\int_E g(x,\mu)\,\mathrm{d}\mu(x)\geqslant\mu(E)^2g(0).\label{align:ch03-3.23}
\end{align}
\end{lemma}

\begin{proof}
不妨设 $E$ 为紧集。因 $g$ 径向递减，对任意 $x,y\in E$ 有 $g(x-y)\geqslant g(\operatorname{diam}E)$。故
\begin{align*}
\int_E g(x,\mu)\,\mathrm{d}\mu(x)
=\int_E\int_E g(x-y)\,\mathrm{d}\mu(y)\,\mathrm{d}\mu(x)
\geqslant g(\operatorname{diam}E)\mu(E)^2.
\end{align*}
若 $\operatorname{diam}E=0$，即 $E$ 为单点，则 $g(0)$ 替换 $g(\operatorname{diam}E)$，结论成立。
\end{proof}

\subsection{容量 $C_{g,p}$}\label{subsection:ch03-容量}

设 $g(x,y)\equiv g(x-y)$，$g:\mathbb{R}^N\to[0,\infty)$ 为径向递减连续函数。

\begin{definition}\label{definition:ch03-Cgp}
设 $p\in(1,\infty)$。对任意 $E\subseteq\mathbb{R}^N$，$E$ 的 $(g,p)$-容量定义为
\begin{align}
C_{g,p}(E)\triangleq \inf_{f\in V_{E,p}}\|f\|_p^p,\qquad
V_{E,p}\triangleq \{f\in L^p(\mathbb{R}^N)\mid (g*f)(x)\geqslant1\ \forall x\in E\}.\label{align:ch03-3.24a}
\end{align}
若 $V_{E,p}=\varnothing $，则 $C_{g,p}(E)\triangleq \infty$。

$E$ 的 $(g,1)$-容量定义为
\begin{align}
C_{g,1}(E)\triangleq \inf_{\sigma\in V_{E,1}}\|\sigma\|,\label{align:ch03-3.24b}
\end{align}
其中
\begin{align}
V_{E,1}\triangleq \left\{\sigma\in\mathscr{R}_f(\mathbb{R}^N)\ \middle|\ \int_{\mathbb{R}^N}g(x-y)\,\mathrm{d}\sigma(y)\geqslant1\ \forall x\in E\right\}.\label{align:ch03-3.24c}
\end{align}
若 $V_{E,1}=\varnothing $，则 $C_{g,1}(E)\triangleq \infty$。
\end{definition}

\begin{remark}\label{remark:ch03-非负可去}
在\eqref{align:ch03-3.24a}中显然可限制 $f\geqslant0$；类似地在\eqref{align:ch03-3.24b}中可限制 $\sigma\geqslant0$。
\end{remark}

\begin{theorem}\label{theorem:ch03-Cgp外容量}
设 $p\in[1,\infty)$。则 $C_{g,p}$ 是 $\mathbb{R}^N$ 上的外容量。
\end{theorem}

\begin{proof}
显然 $C_{g,p}(\varnothing )=0$，且若 $E_1\subseteq E_2$ 则 $C_{g,p}(E_1)\leqslant C_{g,p}(E_2)$。次可列可加性可由标准对角论证得到。右连续性类似得到。
\end{proof}

\begin{proposition}\label{proposition:ch03-Borel可容}
设 $p\in(1,\infty)$。$\mathbb{R}^N$ 的每个 Borel 子集为 $C_{g,p}$-可容集。
\end{proposition}

\begin{proposition}\label{proposition:ch03-对偶}
设 $p\in(1,\infty)$，$q=p/(p-1)$。对任意紧集 $K\subseteq\mathbb{R}^N$，
\begin{align}
C_{g,p}(K)=\sup\left\{\mu(K)^q\ \middle|\ \mu\in\mathscr{R}^+(K),\ \int_{\mathbb{R}^N}g(x-y)\,\mathrm{d}\mu(y)\,\mathrm{d}\mu(x)\leqslant1\right\}^{1/(q-1)}.\label{align:ch03-3.25}
\end{align}
\end{proposition}

\begin{proof}
这是经典对偶定理：在约束 $g*f\geqslant1$（在 $K$ 上）下极小化 $\|f\|_p^p$，等价于在能量约束 $\int g*\mathrm{d}\mu\,\mathrm{d}\mu\leqslant1$ 下极大化 $\mu(K)$。
\end{proof}

\begin{proposition}\label{proposition:ch03-比较}
设 $g_1,g_2:\mathbb{R}^N\to[0,\infty)$ 为径向递减连续函数，$g_1\leqslant g_2$。则对任意 $E\subseteq\mathbb{R}^N$ 与 $p\in[1,\infty)$，
\begin{align*}
C_{g_1,p}(E)\leqslant C_{g_2,p}(E).
\end{align*}
\end{proposition}

\begin{proposition}\label{proposition:ch03-简单函数逼近}
设 $p\in(1,\infty)$，$g\in L^{q/(q-1)}(\mathbb{R}^N)$。则对任意紧集 $K\subseteq\mathbb{R}^N$，映射 $p\mapsto C_{g,p}(K)$ 在 $[1,\infty)$ 上右连续。
\end{proposition}

\begin{lemma}\label{lemma:ch03-位势下半连续}
设 $f\in L^p(\mathbb{R}^N)$（$p\in[1,\infty)$）。则 $g*f$ 下半连续。
\end{lemma}

\begin{proof}
若 $f\geqslant0$，则 $g*f(x)=\int g(x-y)f(y)\,\mathrm{d}y$ 为连续函数的上包络，故下半连续。一般情形由 $f=f^+-f^-$ 分解得到。
\end{proof}

\subsection{$C_{g,p}$ 对 $p$ 的依赖}\label{subsection:ch03-p依赖}

设 $g(x,y)\equiv g(x-y)$。假设
\begin{align}
g(0)<\infty.\label{align:ch03-3.35}
\end{align}
此外假设存在 $\alpha>1$ 使
\begin{align}
\int_{\mathbb{R}^N}g(x)^\alpha\,\mathrm{d}x<\infty.\label{align:ch03-3.40}
\end{align}

\begin{lemma}\label{lemma:ch03-p趋近1}
设 $p\in(1,\infty)$，\eqref{align:ch03-3.35}成立。则对任意紧集 $K\subseteq\mathbb{R}^N$，
\begin{align*}
\lim_{p\to1^+}C_{g,p}(K)=C_{g,1}(K).
\end{align*}
\end{lemma}

\begin{proposition}\label{proposition:ch03-p趋近1}
设 $p\in(1,\infty)$，\eqref{align:ch03-3.35}与\eqref{align:ch03-3.40}成立。则对任意紧集 $K\subseteq\mathbb{R}^N$，映射 $p\mapsto C_{g,p}(K)$ 在 $[1,\alpha')$ 上右连续，其中 $\alpha'>1$ 由 $\frac1\alpha+\frac1{\alpha'}=1$ 定义。
\end{proposition}

\begin{proof}
由\eqref{align:ch03-3.40}，$g\in L^\alpha(\mathbb{R}^N)$。若 $p<\alpha'$，则 $q=p/(p-1)>\alpha$，故 $g\in L^q(\mathbb{R}^N)$。由\cref{proposition:ch03-简单函数逼近}，结论成立。
\end{proof}

\begin{remark}\label{remark:ch03-条件替换}
若将\eqref{align:ch03-3.40}替换为较弱的条件
\begin{align}
\int_{\mathbb{R}^N}g(x)^{\alpha}\,\mathrm{d}x<\infty\label{align:ch03-3.40b}
\end{align}
对某个 $\alpha>N/(N-r)$（$r>0$），则结论在 $[1,\alpha')$ 上成立。
\end{remark}




\end{document}